\section{Results \& Discussion}

% ------------------------------------------------------------------
\subsection{Quantitative Results}

\subsubsection{Distributional Alignment (FID \& KID)}
A well-known limitation of FID is its high bias under finite sample sizes, especially problematic in medical imaging. We address this by extrapolating FID linearly from subset fractions \{25\%, 50\%, 75\%, 100\%\} to an infinite sample limit ($1/N \to 0$). Figure~\ref{fig:fid_curve} illustrates this extrapolation. Mask2Derm achieves a \textbf{bias-corrected FID of 62.33}, competitive with StyleGAN2-ADA on the same HAM10000 corpus ($\sim$65.5) and substantially better than unguided Stable Diffusion (145.26). The Kernel Inception Distance is \textbf{KID $= 0.07 \pm 0.01$}, confirming distributional alignment on this unbiased estimator.

\begin{figure}[H]
    \centering
    \begin{subfigure}[b]{0.48\textwidth}
        \centering
        \safeimg[width=\linewidth]{figures/fid_extrapolation_curve.png}
        \caption{FID Extrapolation}
        \label{fig:fid_curve}
    \end{subfigure}
    \hfill
    \begin{subfigure}[b]{0.48\textwidth}
        \centering
        \safeimg[width=\linewidth]{figures/iou_distribution.png}
        \caption{Shape Consistency (IoU)}
        \label{fig:iou_dist}
    \end{subfigure}
    \caption{Quantitative evaluation: (a) FID extrapolation to infinite sample size; (b) distribution of IoU scores measuring mask adherence.}
    \label{fig:quantitative_metrics}
\end{figure}

\subsubsection{Geometric Fidelity (Shape Consistency)}
We generated 5,312 synthetic images using HAM10000 masks and processed them through a pre-trained DeepLabV3+ (ResNet101 backbone). The model achieved \textbf{mean IoU $= 0.8658 \pm 0.1187$} and \textbf{mean DSC $= 0.9222 \pm 0.0955$}. Figure~\ref{fig:iou_dist} shows a strongly left-skewed IoU distribution with the majority of samples exceeding 0.90, confirming that ControlNet effectively suppresses boundary hallucination common in unconditioned generative models.

% ------------------------------------------------------------------
\subsection{Qualitative Results}

Figure~\ref{fig:qualitative} presents real HAM10000 images (top row) alongside synthetic images generated from their corresponding masks (bottom row). The generated lesions reproduce complex stochastic textures—variegated pigmentation, irregular borders, pigment networks, dots and globules—at a quality visually indistinguishable from real acquisitions. The optics simulation endows synthetic images with the characteristic circular field-of-view, radial light falloff, and barrel distortion of real dermoscopes, preventing the ``digitally perfect'' look that typically betrays synthetic data.

\begin{figure}[t!]
    \centering
    \safeimg[width=\linewidth]{figures/qualitative_comparison.png}
    \caption{Qualitative comparison: real images from HAM10000 with lens simulation (top) versus synthetic images from Mask2Derm (bottom). Texture, color, and optical characteristics are closely matched.}
    \label{fig:qualitative}
\end{figure}

Figure~\ref{fig:benign_malignant} demonstrates class-conditional generation. Malignant samples (top two rows) exhibit asymmetric borders, color variegation, and complex internal structures, while benign samples (bottom two rows) display uniform pigmentation and regular symmetrical shapes consistent with clinical presentations of nevi.

\begin{figure}[t!]
    \centering
    \safeimg[width=\linewidth]{figures/benign_malignant_comparison.png}
    \caption{Class-conditional generation: malignant lesions (top two rows) versus benign lesions (bottom two rows). The model captures class-discriminative morphological features.}
    \label{fig:benign_malignant}
\end{figure}

% ------------------------------------------------------------------
\subsection{Train on Synthetic, Test on Real (TSTR)}

\textbf{Setup.} We fine-tuned a DeepLabV3+ segmentation model ($M_\text{synth}$) using \emph{only} the 5,312 synthetic image–mask pairs from Mask2Derm. As a baseline ($M_\text{base}$), we evaluated a DeepLabV3+ model initialized with ImageNet weights and not exposed to any dermoscopy-specific data. Both models were evaluated on the held-out real HAM10000 test set.

\textbf{Results.} $M_\text{synth}$ achieved mean IoU of \textbf{$0.9254 \pm 0.0412$} and mean DSC of \textbf{0.9608}, significantly outperforming $M_\text{base}$ (IoU $= 0.7910 \pm 0.1625$, DSC $= 0.8716$). The performance delta $\Delta = +0.1344$ IoU confirms that Mask2Derm generates semantically and structurally accurate data that generalizes zero-shot to real patient images.

\begin{figure}[H]
    \centering
    \safeimg[width=\linewidth]{figures/tstr_iou_boxplot.png}
    \caption{TSTR comparison: IoU distribution of the synthetic-finetuned model versus the ImageNet baseline on the real test set. The synthetic model exhibits superior median performance and lower variance.}
    \label{fig:tstr_boxplot}
\end{figure}

\textbf{Interpretation.} The relatively low performance of the ImageNet baseline highlights the significant domain gap between natural images and dermoscopy. Our synthetic data successfully bridges this gap: by training exclusively on generated proxies, the model acquires domain-specific priors for lesion boundaries, texture gradients, and skin contrast that transfer to real data. This confirms that Mask2Derm is not merely visually coherent but carries genuine clinical signal.

% ------------------------------------------------------------------
\subsection{Limitations}

Despite strong results, several limitations warrant explicit discussion. First, Mask2Derm primarily generates lesions on fair skin tones, reflecting the bias jointly inherited from the underlying Stable Diffusion model and the HAM10000 dataset, which is heavily skewed toward European cohorts. This skin-tone bias is a critical concern for equitable clinical deployment and must be addressed before real-world adoption. Second, the physics-based optics simulation is a first-order approximation; it does not model polarized light, subsurface scattering, or device-specific chromatic aberrations. Third, the current FID of 62.33, while competitive on the full heterogeneous HAM10000 corpus, remains above the scores achieved by specialized models trained on curated single-class subsets. Finally, our TSTR evaluation uses a single segmentation architecture (DeepLabV3+); broader downstream validation across multiple architectures and tasks would further strengthen the clinical utility claim.
